% REVIEW DRAFT: newly written from the workshop proposal and public information.
% Confirm the text and organizer signatures before publication.
Welcome to the Workshop on Multi-modality and Multi-lingual Language Learning
(Multi-LLL 2026), held on November 9, 2026, in Hengqin, Zhuhai, China,
in conjunction with AACL-IJCNLP 2026.

Language learning draws on speech, writing, visual information, and interaction.
This workshop brings together researchers working across natural language
processing, computer vision, speech, artificial intelligence, and education to
explore how these resources can support learners with different linguistic
backgrounds and needs. Its scope includes learning technologies, assessment,
feedback, and applications for English, Chinese, and low-resource languages.

Alongside research presentations and invited talks, the workshop features the
VisCGEC-v2.0 shared task on visual Chinese grammatical error correction.
The task asks systems to locate and correct errors in images of handwritten
Chinese text, with sentence-level and paragraph-level tracks. It connects
visual understanding with language analysis in an educational setting.

We thank the authors, reviewers, invited speakers, shared-task participants,
and the AACL-IJCNLP 2026 organizing team for their contributions. We hope the
workshop encourages continued exchange and collaboration on language learning
across modalities and languages.

\bigskip
\noindent Xiaoman Wang, Qianyu Wang, Zhengyuan Liu, Nancy F. Chen,\\
Zheng Yuan, Yunshi Lan, and Jing Jiang\\
Multi-LLL 2026 Workshop Organizers
